\documentclass[10pt,a4paper]{article}
\usepackage{amsmath,amssymb,amsthm,graphicx}
\usepackage{amsmath,amssymb,amsthm,mathtools}
\usepackage[no-math]{fontspec}
\newenvironment{authorclaim}[1]{\par\medskip\noindent\textbf{#1.}\ \ignorespaces}{\par\medskip}
\allowdisplaybreaks
\title{Existence of palette-alternating tree codes with Gilbert-Varshamov distance-rate tradeoff}
\author{Gil Cohen, Shahar Samocha}
\date{doi:10.4230/LIPIcs.CCC.2020.11; CC BY 3.0 Unported}
\begin{document}
\maketitle
\noindent\textit{Editorial scope.} One counted current-formal palette-alternating existence/GV mechanism. Four-color common construction/support and interactive application are context/support. The actual main probability proof uses union bounds/entropy/geometric sums;no invented LLL proof. Historical TR19-147 withdrawal concerns an already-achieved novelty claim;no first-capacity claim or old-version proof used.\par\medskip
\noindent\textit{Allowed standard prerequisites (editorial).}\begin{itemize}
\item Finite fields F2/F4,linear trace maps,binary entropy bounds and probabilistic union bounds. Formal Lemmas7–8 retained as cited standard inequalities;complete authors' field construction,independence/common-core and all union/entropy/geometric-tail arguments are retained. Full32 original bibliography entries retained.
\end{itemize}\medskip
\section*{Tree codes and palette-alternating models}
We turn to formally define tree codes. Let $T$ be a rooted binary tree that is endowed with an edge coloring from some ambient color set (or alphabet) $\Sigma$. Let $u,v$ be a pair of vertices in $T$ with equal depth and a least common ancestor $w$. Let $\ell$ be the distance, in edges, from $u$ to $w$. Let $p_u,p_v\in\Sigma^\ell$ be the sequences of colors on the path from $w$ to $u$ and to $v$, respectively. We define $h(u,v)$ to be the relative Hamming distance between $p_u$ and $p_v$.
\begin{authorclaim}{Definition1 (Tree codes [29])}
Let $T$ be the complete rooted infinite binary tree. The tree $T$, together with an edge-coloring of $T$ by a color set $\Sigma$ is called a tree code with distance $\delta$ if for every pair of vertices $u,v$ with equal depth it holds that $h(u,v)\ge\delta$. When there is no $\delta>0$ for which $T$ is a tree code with distance $\delta$ we say that $T$ has vanishing distance.
\end{authorclaim}
Schulman [29] proved that for every distance parameter $\delta<1$ tree codes with a constant number of colors $c=c(\delta)$ exist. Although tree codes are used in different ways by different interactive coding schemes, one aspect is common to all: When a party wishes to send a bit, a suitable color from $\Sigma$ is sent instead. Thus, the rate of all tree-code-based coding schemes is bounded above by $1/\log_2|\Sigma|$. One is led to ask a natural combinatorial question--what is the least number of colors in a tree code with non-vanishing distance?
\paragraph{1.1 Tree codes:4 colors suffice and are necessary}
We first observe that 3 colors do not suffice and, as a result, the rate of every tree-code-based coding scheme cannot exceed1/2, let alone approach capacity. Consider any 3-color tree code. First, we may assume that every two siblings are connected to their parent with edges having distinct colors as otherwise the distance of the tree code is 0. Let $u,v$ be any two vertices. Out of $u,v$ go four edges and so by the pigeonhole principal in every 3-coloring, two of these edges share the same color. By the above, one of these edges goes out of $u$ and the other goes out of $v$. This implies that, starting from the two sons of the root, one can construct two paths of any desired length $n\ge1$ with the same color pattern, establishing that the tree has vanishing distance.
Based on the ideas Schulman introduced to prove the existence of tree codes with a constant number of colors, we complement the above observation and establish that 4 colors suffice for a tree code with non-vanishing distance.
\begin{authorclaim}{Theorem2}
There exists a 4-color tree code with distance 0.136.
\end{authorclaim}
\paragraph{Definition3 (Palette-alternating tree codes)}
Let $\Sigma_0,\ldots,\Sigma_{c-1}$ be (not necessarily distinct) sets. Let $T$ be the complete rooted infinite binary tree. A palette-alternating tree code is an edge-coloring of $T$ where at layer $t\in\mathbb N$ the colors are taken from the set $\Sigma_{t\bmod c}$. $T$ is said to have distance $\delta$ if for every pair of vertices $u,v$ with equal depth it holds that $h(u,v)\ge\delta$. We define the rate $\rho$ of $T$ to be the number satisfying
\[
\frac1\rho=\frac1c\sum_{i=0}^{c-1}\log_2|\Sigma_i|.
\]
We suggest that the flexibility introduced by palette-alternating tree codes allows one to better capture the notion of rate in the online setting. Indeed, the importance of rate is only significant when "long" messages are being sent and so, informally, using a big palette of colors only once in a while should not be considered as an indication of poor rate. Our definition of rate formalizes that property. Note that we still insist on having the distance measured in terms of worst-case--a must as we wish to replace tree codes with palette-alternating tree codes in interactive coding schemes. It is only the rate that is being, in a sense, amortized.
As mentioned, we prove that palette-alternating tree codes can have rate approaching arbitrarily close to1 while maintaining non-vanishing distance, thus bypass the1/2 bound proven for (standard) tree codes.
\section*{2 Preliminaries}
Unless otherwise stated, all logarithms are taken to the base2. We denote by $\mathbb N$ the set of natural numbers (of course, including0), and write $\mathbb N_1$ for $\mathbb N\setminus\{0\}$. For integers $a\le b$ we write $[a,b]$ for all integers in this interval. For an integer $c\ge1$, we let $[c]=\{1,2,\ldots,c\}$. We follow the convention that strings are indexed starting from1. For two strings $x,y\in\Sigma_1\times\cdots\times\Sigma_n$, we denote by $\Delta(x,y)$ their hamming distance. We make use of the following standard inequalities.
\begin{authorclaim}{Lemma7}
For every integers $1\le k\le n$ with $k/n=\delta\le1/2$ it holds that $\sum_{i=0}^k\binom ni\le2^{H(\delta)n}$.
\end{authorclaim}
\begin{authorclaim}{Lemma8}
For every $0<x<1/2$ it holds that $x/(2\log_2(6/x))\le H^{-1}(x)\le x/\log_2(1/x)$.
\end{authorclaim}

\begin{authorclaim}{Theorem4}
For every $\epsilon>0$ there exists a palette-alternating tree code with rate $1-\epsilon$ and distance $\delta=\Omega(\epsilon\cdot\log^{-1}(1/\epsilon))$.
\end{authorclaim}
\paragraph{Comparison with the Gilbert-Varshamov bound}
Observe that the distance-rate trade-off obtained in Theorem4 is the same as the one obtained by the Gilbert-Varshamov bound for standard offline binary error correcting codes, and in particular is optimal (up to constant factors). Interestingly, while it is known that the channel capacity in the online setting is $1-\Theta(\sqrt{\epsilon\log(1/\epsilon)})$--significantly lower than in the offline setting [20], the online requirement on the encoding function itself does not cost more in terms of the distance-rate trade-off. Rather, it is the additional overhead incurred by the mechanism required for synchronization that is responsible for the lower channel capacity in the online setting. We elaborate more on this in Section1.3.
The proof of Theorem4, which can be found in Section4, is based on a variant of the construction we use in Theorem2. There, the alphabet symbols are taken from the field of four elements, $\mathbb F_4$. The key idea in obtaining the savings in the alphabet size is to trace the $\mathbb F_4$ field elements down to $\mathbb F_2$ in most layers. Interestingly, we cannot afford to work over the field $\mathbb F_3$ as we crucially rely on the fact that the characteristic of the fields is2 as well as on the smaller field being a subfield of the larger one.

\section*{3 Binary Tree Codes:Four Colors Suffice}
In this section we prove Theorem2. We start by setting some notation. Let $T$ be the infinite complete rooted binary tree. We identify length-$n$ paths in $T$ that starts at the root with length-$n$ binary strings in the natural way. Namely, we identify left son and right son with 0 and 1, respectively. Given a node $v$ at depth $n\ge1$ we define $p_v\in\{0,1\}^n$ to be the string that encodes the (unique) path from the root to $v$.
An edge-coloring of $T$ by a color set $\Sigma$ is given by a function, which for ease of readability, we slightly abuse notation and also denote by $T:\{0,1\}^{\mathbb N_1}\to\Sigma^{\mathbb N_1}$, where the color of an edge $e=u\to v$ is $T(p_v)_{\operatorname{depth}(v)}$. Note that $T$ is an online function, namely, for every $x\in\{0,1\}^{\mathbb N_1}$ and $i\in\mathbb N_1$, the value $T(x)_i$ is determined by $x_1,\ldots,x_i$.
\paragraph{The (probabilistic) construction}
Let $\{R_i\}_{i\in\mathbb N_1}$ be a sequence of independent random variables, each is uniformly distributed over $\mathbb F_4$--the field of 4 elements. Let $\mathbb F_2$ be the (unique) subfield of $\mathbb F_4$ of size 2. Define the (random) coloring function $T:\mathbb F_2^{\mathbb N_1}\to\mathbb F_4^{\mathbb N_1}$ (where we identify $\mathbb F_2$ and $\{0,1\}$ in the natural way) as follows: for every $t\in\mathbb N_1$
\begin{equation}
T(x)_t=\sum_{i=1}^tR_{t+1-i}x_i.\tag{1}
\end{equation}
\begin{authorclaim}{Definition9}
Let $v$ be a depth-$n$ vertex in $T$. Let $\ell\ge1$ and $x,y\in\mathbb F_2^{\ell-1}$. For $k=1,\ldots,\ell$ we define the random variable
\[
a_v(x,y,k)=T(p_v\circ1\circ y)_{n+k}-T(p_v\circ0\circ x)_{n+k}.
\]
Note that $a_v(x,y,k)$ is a (random) element in $\mathbb F_4$. We define the integral random variable $h_v(x,y)=\sum_{k=1}^\ell I_k$, where $I_k$ is the indicator random variable that equals 1 when $a_v(x,y,k)\ne0$. Note that $h_v(x,y)\in\{0,1,\ldots,\ell\}$ is the Hamming distance between $T(p_v\circ0\circ x)_{[n+1,n+\ell]}$ and $T(p_v\circ1\circ y)_{[n+1,n+\ell]}$.
\end{authorclaim}
\begin{authorclaim}{Claim10}
Let $v$ be a vertex in $T$. Let $\ell\ge1$ and $x,y\in\mathbb F_2^{\ell-1}$. Then, for every $k\in\{1,\ldots,\ell\}$ it holds that $a_v(x,y,k)=R_k+\sum_{i=1}^{k-1}R_{k-i}(y-x)_i$.
\end{authorclaim}
\begin{proof}
Denote the depth of $v$ by $n$. Fix $k\in\{1,\ldots,\ell\}$. By Equation(1),
\begin{align*}
T(p_v\circ0\circ x)_{n+k}
&=\sum_{i=1}^{n+k}R_{n+k+1-i}(p_v\circ0\circ x)_i\\
&=\sum_{i=1}^nR_{n+k+1-i}(p_v)_i+\sum_{i=1}^kR_{k+1-i}(0\circ x)_i.
\end{align*}
Similarly
\[
T(p_v\circ1\circ y)_{n+k}
=\sum_{i=1}^nR_{n+k+1-i}(p_v)_i+\sum_{i=1}^kR_{k+1-i}(1\circ y)_i.
\]
Thus,
\begin{align*}
a_v(x,y,k)
&=\sum_{i=1}^kR_{k+1-i}(1\circ y)_i-\sum_{i=1}^kR_{k+1-i}(0\circ x)_i\\
&=R_k+\sum_{i=1}^{k-1}R_{k-i}(y-x)_i.
\end{align*}
\end{proof}
\begin{authorclaim}{Claim11}
Let $v$ be a vertex in $T$. Let $\ell\ge1$ and $x,y\in\mathbb F_2^{\ell-1}$. Then, the random variables $a_v(x,y,1),\ldots,a_v(x,y,\ell)$ are independent and each is uniformly distributed over $\mathbb F_4$.
\end{authorclaim}
\begin{proof}
By Claim10, $a_v(x,y,k)=R_k+L_k$ where $L_k$ is some $\mathbb F_4$-linear combination of $R_1,\ldots,R_{k-1}$. Therefore, $a_v(x,y,k)$ is independent of the joint distribution of $a_v(x,y,1),\ldots,a_v(x,y,k-1)$. As this holds for every $k$ we have that $a_v(x,y,1),\ldots,a_v(x,y,\ell)$ are independent. To conclude the proof, note that for every fixing of $R_1,\ldots,R_{k-1}$, $a_v(x,y,k)=R_k+\ell_k$ for some fixed $\ell_k\in\mathbb F_4$ and so $a_v(x,y,k)$ is uniform over $\mathbb F_4$.
\end{proof}
\begin{authorclaim}{Claim12}
For every two vertices $u,v$ in $T$ and every $x,y\in\mathbb F_2^{\ell-1}$, $h_v(x,y)=h_u(x,y)$, and $h_v(x,y)=h_v(0^{\ell-1},y-x)$.
\end{authorclaim}
\begin{proof}
The first equality follows immediately by Claim10 as, for every $k\in\{1,\ldots,\ell\}$, the expression obtained for $a_v(x,y,k)$ is independent of the choice of $v$. As for the second asserted equality, again by Claim10,
\begin{align*}
a_v(x,y,k)&=R_k+\sum_{i=1}^{k-1}R_{k-i}(y-x)_i\\
&=R_k+\sum_{i=1}^{k-1}R_{k-i}((y-x)-0)_i
=a_v(0^{\ell-1},y-x,k),
\end{align*}
where observe that for the last equality we are using the fact that $\mathbb F_2$ is a subfield of $\mathbb F_4$ and so $y-x\in\mathbb F_2^{\ell-1}$. Indeed, recall that $a_v$'s second argument is a binary string and so the equality above would have been meaningless otherwise. The above equation implies $h_v(x,y)=h_v(0^{\ell-1},y-x)$, proving the claim.
\end{proof}
Given Claim12 we can simplify our notation as follows. Let $r$ denote the root of $T$. For $x\in\{0,1\}^{\ell-1}$ and $k\in\{1,\ldots,\ell\}$ we define the random variables $a(x,k)=a_r(0^\ell,1\circ x,k)$, $h(x)=h_r(0^{\ell-1},x)$. Note that $h(x)=\sum_{k=1}^\ell a(x,k)$.
\begin{authorclaim}{Theorem13}
There exists a fixing of the sequence $\{R_i\}_i$ such that the function $T$ is a tree code with distance 0.05.
\end{authorclaim}
\begin{proof}
First note that for every fixing of the sequence $\{R_i\}_i$, $T$ is an online function. Observe that, for a fixing of $\{R_i\}_i$, $T$ is a tree code with distance $\delta$ if and only if for every $\ell\ge1$ and $x\in\{0,1\}^{\ell-1}$ it holds that $h(x)\ge\delta\ell$. Indeed, recall that by definition, $T$ is a tree code with distance $\delta$ if and only if for every vertex $v$ in $T$, $\ell\ge1$, and for every $x,y\in\{0,1\}^{\ell-1}$ it holds that $h_v(x,y)\ge\delta\ell$. However, by Claim12, $h_v(x,y)=h(y-x)$.
For $x\in\{0,1\}^{\ell-1}$ denote by $E(x)$ the event $h(x)<\delta\ell$. By the above discussion, it suffices to prove, for $\delta=0.05$, that $\Pr[\bigcup_{x\in\{0,1\}^{\mathbb N}}E(x)]<1$. To this end, by the union bound, it suffices to prove that $\sum_{x\in\{0,1\}^{\mathbb N}}\Pr[E(x)]<1$.
Consider any $x\in\{0,1\}^{\ell-1}$ with $\ell\ge1$. Note that the event $E(x)$ holds if and only if there exists a set $\mathcal T\subseteq\{1,\ldots,\ell\}$ of size $|\mathcal T|\ge\lceil(1-\delta)\ell\rceil$ such that for every $k\in\mathcal T$, $a(x,k)=0$. By taking the union bound over all such sets $\mathcal T$, and using that $a(x,1),\ldots,a(x,\ell)$ are independent and each is uniformly distributed over $\mathbb F_4$ (Claim11), we get
\begin{equation}
\Pr[E(x)]\le\binom\ell{\lceil(1-\delta)\ell\rceil}4^{-\lceil(1-\delta)\ell\rceil}.\tag{2}
\end{equation}
By Lemma7, we have that
\[
\frac1\ell\log_2\binom\ell{\lceil(1-\delta)\ell\rceil}
\le H\left(\frac{\lceil(1-\delta)\ell\rceil}{\ell}\right).
\]
As $\delta<1/2$ and since the entropy function $H$ decreases in $[1/2,1]$ we have that
$H(\lceil(1-\delta)\ell\rceil/\ell)\le H(1-\delta)=H(\delta)$.
Substitute to Equation(2), we get that $\Pr[E(x)]\le2^{(H(\delta)-2(1-\delta))\ell}$. Thus,
\begin{align*}
\sum_{x\in\{0,1\}^{\mathbb N}}\Pr[E(x)]
&\le\sum_{\ell=1}^\infty2^{\ell-1}\cdot2^{(H(\delta)-2(1-\delta))\ell}\\
&=\frac12\sum_{\ell=1}^\infty2^{(H(\delta)+2\delta-1)\ell}.
\end{align*}
One can verify that for $\delta=0.05$ the above geometric sum is strictly smaller than1, and the theorem follows.
\end{proof}
\subsection*{3.1 Improving the distance}
We now show a method for improving the distance. We illustrate it to obtain a bound of 0.136 on the distance, which proves Theorem2, though we believe that the method can be used to push the bound further. It is fairly easy to show that the distance of a 4-color tree code cannot be larger than 1/2.
\begin{authorclaim}{Theorem14}
There exists a fixing of the sequence $\{R_i\}_i$ such that the function $T$ is a tree code with distance 0.136.
\end{authorclaim}
\begin{proof}
For the proof it will be convenient to consider a specific representation of $\mathbb F_4$. We make use of the standard construction of $\mathbb F_4$ as a quotient of the polynomial ring over $\mathbb F_2$ with respect to an ideal generated by a degree 2 irreducible element as follows. Note that $t^2+t+1\in\mathbb F_2[t]$ is irreducible, and so $K=\mathbb F_2[t]/\langle t^2+t+1\rangle$ is a field of 4 elements which we will take as the construction for $\mathbb F_4$. Let $\alpha$ be the class of $t$ in $K$. In this representation, the field $\mathbb F_4$ consists of the elements $0,1,\alpha,\alpha+1$ where $\alpha^2+\alpha+1=0$.
Consider the sequence $\{R_i\}_{i\in\mathbb N}$ as in the beginning of the section but with the fixings $R_1=1$ and $R_2=\alpha$. Observe that for every $x\in\mathbb F_2^{\ell-1}$ with $\ell\ge2$ it holds that $a(x,1)=1$ and $a(x,2)=\alpha+x_1$. In particular, $a(x,1),a(x,2)$ are both non-zeros and so $h(x)\ge2$. Let $\ell_0\ge2$ be an integer parameter to be chosen later on. By the above, we have that for every $x\in\mathbb F_2^{\ell-1}$ with $\ell\le\ell_0$ it holds that
\begin{equation}
\frac{h(x)}{\ell}\ge\frac2{\ell_0}.\tag{3}
\end{equation}
For $x\in\{0,1\}^{\ell-1}$ denote by $E_{1,\alpha}(x)$ the event $h(x)<\delta\ell$ with the $\{R_i\}_{i\in\mathbb N}$ as defined above, namely, $R_1=1$, $R_2=\alpha$ and the rest of the random variables $\{R_i\mid i\ge3\}$ are independent and uniformly distributed over $\mathbb F_4$. Once we establish a bound of
\begin{equation}
\Pr[\bigcup_{|x|\ge\ell_0}E_{1,\alpha}(x)]<1\tag{4}
\end{equation}
for some choice of $\delta$ then, combined with Equation(3), we will establish the existence of a tree code with distance at least $\min(2/\ell_0,\delta)$.
Consider any $x\in\{0,1\}^{\ell-1}$ with $\ell\ge\ell_0+1$. The event $E_{1,\alpha}(x)$ holds if and only if there exists a set $\mathcal T\subseteq\{3,\ldots,\ell\}$ of size $|\mathcal T|\ge\lceil(1-\delta)\ell\rceil$ such that for every $k\in\mathcal T$, $a(x,k)=0$. By taking the union bound over all such sets $\mathcal T$, and using that $a(x,3),\ldots,a(x,\ell)$ are independent and each is uniformly distributed over $\mathbb F_4$, we get that
\begin{align*}
\Pr[E_{1,\alpha}(x)]
&\le\binom{\ell-2}{\lceil(1-\delta)\ell\rceil}4^{-\lceil(1-\delta)\ell\rceil}\\
&\le\binom\ell{\lceil(1-\delta)\ell\rceil}4^{-\lceil(1-\delta)\ell\rceil}.
\end{align*}
By Lemma7, we have that
\[
\frac1\ell\log_2\binom\ell{\lceil(1-\delta)\ell\rceil}
\le H\left(\frac{\lceil(1-\delta)\ell\rceil}{\ell}\right).
\]
As we will choose $\delta<1/2$ and the entropy function $H$ decreases in $[1/2,1]$ we have that $H(\lceil(1-\delta)\ell\rceil/\ell)\le H(1-\delta)=H(\delta)$.
Thus, $\Pr[E_{1,\alpha}(x)]\le2^{(H(\delta)-2(1-\delta))\ell}$. By substituting the above equation to Equation(4), we get that
\[
\sum_{|x|\ge\ell_0}\Pr[E_{1,\alpha}(x)]
\le\sum_{\ell=\ell_0+1}^\infty2^{\ell-1}\cdot2^{(H(\delta)-2(1-\delta))\ell}.
\]
Write $\beta=2^{H(\delta)+2\delta-1}$. Then, the above is bounded by $\frac12\sum_{\ell=\ell_0+1}^\infty\beta^\ell=\beta^{\ell_0+1}/(2(1-\beta))$.
Consider the real polynomial $f_{\ell_0}(x)=x^{\ell_0+1}-2(1-x)$. We have that $f'_{\ell_0}(x)=(\ell_0+1)x^{\ell_0}+2$.
Since $\ell_0\ge2$, $f'_{\ell_0}(x)>0$ for all $x\ge0$. Further, $f_{\ell_0}(0)=-2$ and $f_{\ell_0}(1)=1$. Thus, $f_{\ell_0}(x)$ has a single root $\beta_{\ell_0}\in[0,1]$ (in fact, $\beta_{\ell_0}$ is monotone-increasing as a function of $\ell_0$, and $\beta_{\ell_0}\to1$ as $\ell_0\to\infty$). For a fixed choice of $\ell_0$, by choosing $\beta<\beta_{\ell_0}$ and solving for $\delta$ (recall $\beta=2^{H(\delta)+2\delta-1}$) to obtain $\delta_{\ell_0}$, we get that there exists a fixing of $\{R_i\mid i\ge3\}$ such that the obtained tree code has distance at least $\min(\delta_{\ell_0},2/\ell_0)$. Thus, the obtained bound is $\max_{\ell_0\ge2}\min(\delta_{\ell_0},2/\ell_0)$.
One can verify that $\ell_0=14$ maximizes the above equation to get distance larger than 0.136.
\end{proof}

\section*{4 Palette-Alternating Tree Codes}
In this section we prove Theorem4. To this end we recall the definition of the (field) trace function $\operatorname{Tr}:\mathbb F_4\to\mathbb F_2$ that is given by $\operatorname{Tr}(x)=x+x^2$. Observe that the trace function is an $\mathbb F_2$-linear map whose image and kernel are $\mathbb F_2$. In particular, if $X$ is uniform over $\mathbb F_4$, then $\operatorname{Tr}(X)$ is uniform over $\mathbb F_2$.
Let $\epsilon$ be a given parameter and define $b=\lceil1/\epsilon\rceil$. Let $\{R_i\}_{i\in\mathbb N}$ be a sequence of independent random variables, each is uniformly distributed over $\mathbb F_4$ except that $R_1$ is fixed to $R_1=1$. We define a palette-alternating tree code with $b$ palette sets $\Sigma_0,\ldots,\Sigma_{b-1}$ such that $\Sigma_0=\mathbb F_4$ and $\Sigma_i=\mathbb F_2$ for $i>0$. Let $x\in\mathbb F_2^{\mathbb N}$. For every $k\in\mathbb N$, define $S_k(x)=\sum_{i=1}^kR_{k+1-i}x_i$, where addition and multiplication are performed in $\mathbb F_4$ and, as usual, $\mathbb F_2$ is identified with the unique subfield of two elements in $\mathbb F_4$. The coloring function is given by
\[
T(x)_k=\begin{cases}S_k(x),&k\equiv_b0;\\\operatorname{Tr}(S_k(x)),&\text{otherwise}.\end{cases}
\]
\begin{authorclaim}{Theorem15}
The function $T$ above is a palette-alternating tree code with rate $1-\epsilon$ and distance $\delta=\Omega(\epsilon\log^{-1}(1/\epsilon))$.
\end{authorclaim}
\begin{proof}
First, observe that $T$ is indeed an online function with rate larger than $1-\epsilon$. Further Definition9 can be carried over to the more general case of palette-alternating tree codes. We turn to prove an analog to Claim10.
\begin{authorclaim}{Claim16}
Let $v$ be a depth-$n$ vertex in $T$. Let $\ell\ge1$ and $x,y\in\mathbb F_2^{\ell-1}$. Then, for every $k\in\{1,\ldots,\ell\}$ it holds that
\[
a_v(x,y,k)=\begin{cases}R_k+S_{k-1}(y-x),&n+k\equiv_b0;\\
\operatorname{Tr}(R_k+S_{k-1}(y-x)),&\text{otherwise}.\end{cases}
\]
\end{authorclaim}
\begin{proof}
Fix $k\in\{1,\ldots,\ell\}$. Assume first that $n+k\equiv_b0$. Then,
\begin{align*}
T(p_v\circ0\circ x)_{n+k}
&=\sum_{i=1}^{n+k}R_{n+k+1-i}(p_v\circ0\circ x)_i\\
&=\sum_{i=1}^nR_{n+k+1-i}(p_v)_i
+\sum_{i=n+1}^{n+k}R_{n+k+1-i}(0\circ x)_{i-n}\\
&=\sum_{i=1}^nR_{n+k+1-i}(p_v)_i
+\sum_{i=1}^kR_{k+1-i}(0\circ x)_i.
\end{align*}
Similarly
\[
T(p_v\circ1\circ y)_{n+k}
=\sum_{i=1}^nR_{n+k+1-i}(p_v)_i+\sum_{i=1}^kR_{k+1-i}(1\circ y)_i.
\]
Thus,
\begin{align*}
a_v(x,y,k)
&=\sum_{i=1}^kR_{k+1-i}(1\circ y)_i-\sum_{i=1}^kR_{k+1-i}(0\circ x)_i\\
&=R_k+\sum_{i=1}^{k-1}R_{k-i}(y-x)_i=R_k+S_{k-1}(y-x).
\end{align*}
Assume now that $n+k\not\equiv_b0$. Using that $\operatorname{Tr}$ is $\mathbb F_2$-linear,
\begin{align*}
T(p_v\circ0\circ x)_{n+k}
&=\operatorname{Tr}\left(\sum_{i=1}^{n+k}R_{n+k+1-i}(p_v\circ0\circ x)_i\right)\\
&=\operatorname{Tr}\left(\sum_{i=1}^nR_{n+k+1-i}(p_v)_i\right)
+\sum_{i=n+1}^{n+k}\operatorname{Tr}(R_{n+k+1-i})(0\circ x)_{i-n}\\
&=\operatorname{Tr}\left(\sum_{i=1}^nR_{n+k+1-i}(p_v)_i\right)
+\sum_{i=1}^k\operatorname{Tr}(R_{k+1-i})(0\circ x)_i.
\end{align*}
Similarly
\[
T(p_v\circ1\circ y)_{n+k}
=\operatorname{Tr}\left(\sum_{i=1}^nR_{n+k+1-i}(p_v)_i\right)
+\sum_{i=1}^k\operatorname{Tr}(R_{k+1-i})(1\circ y)_i.
\]
Thus, again by $\mathbb F_2$-linearity of $\operatorname{Tr}$,
\[
a_v(x,y,k)=\operatorname{Tr}(R_k)+\sum_{i=1}^{k-1}\operatorname{Tr}(R_{k-i})(y-x)_i
=\operatorname{Tr}(R_k+S_{k-1}(y-x)).
\]
\end{proof}
\begin{authorclaim}{Claim17}
Let $v$ be a depth-$n$ vertex and $x,y\in\mathbb F_2^{\ell-1}$ distinct. Then, the random variables $a_v(x,y,1),\ldots,a_v(x,y,\ell)$ are independent. Moreover, let $k\in[\ell]$. If $n+k\equiv_b0$ then $a_v(x,y,k)$ is uniformly distributed over $\mathbb F_4$ and otherwise it is uniform over $\mathbb F_2$.
\end{authorclaim}
\begin{proof}
By Claim16, if $n+k\equiv_b0$ then $a_v(x,y,k)=R_k+L_k$ where $L_k$ is a linear combination of $R_1,\ldots,R_{k-1}$. Thus, in this case, $a_v(x,y,k)$ is independent of the joint distribution of $a_v(x,y,1),\ldots,a_v(x,y,k-1)$. Otherwise, namely $n+k\not\equiv_b0$, we have that $a_v(x,y,k)=\operatorname{Tr}(R_k+L_k)=\operatorname{Tr}(R_k)+\operatorname{Tr}(L_k)$. Since for every fixing of $L_k$, $a_v(x,y,k)$ is uniform over $\mathbb F_2$, we have that $a_v(x,y,k)$ is independent of the joint distribution of $a_v(x,y,1),\ldots,a_v(x,y,k-1)$. As this holds for every $k\in[\ell]$ we have that $a_v(x,y,1),\ldots,a_v(x,y,\ell)$ are independent and their marginal distributions are as stated.
\end{proof}
\begin{authorclaim}{Claim18}
Let $u,v$ be two vertices with depth $n,m$, respectively such that $n\equiv_bm$. Let $x,y\in\mathbb F_2^{\ell-1}$. Then, $h_v(x,y)=h_u(x,y)$, and $h_v(x,y)=h_v(0^{\ell-1},y-x)$.
\end{authorclaim}
\begin{proof}
Let $C_k=R_k+S_{k-1}(y-x)$. By Claim16,
\[
a_u(x,y,k)=\begin{cases}C_k,&n+k\equiv_b0;\\\operatorname{Tr}(C_k),&\text{otherwise}.\end{cases}
\]
As $C_k$ is independent of the choice of $u$ and $n\equiv_bm$ we have that $a_u(x,y,k)$ is the same random variable as $a_v(x,y,k)$. Since this holds for every $k$, we have that $h_v(x,y)=h_u(x,y)$.
We turn to prove the the second asserted equality. Assume first that $k\in[\ell]$ is such that $n+k\equiv_b0$. By Claim16,
\begin{align*}
a_u(x,y,k)&=R_k+S_{k-1}(y-x)\\
&=R_k+S_{k-1}((y-x)-0^{\ell-1})
=a_u(0^{\ell-1},y-x,k),
\end{align*}
where observe that for the last equality we are using the fact that $\mathbb F_2$ is a subfield of $\mathbb F_4$ and so $y-x\in\mathbb F_2^{\ell-1}$. Indeed, recall that $a_v$'s second argument is a binary string and so the equality above would have been meaningless otherwise. The case $n+k\not\equiv_b0$ follows by a similar argument and using the $\mathbb F_2$-linearity of $\operatorname{Tr}$.
\end{proof}
Given Claim18, we can simplify our notation as follows. Let $v_0$ denote the root of the tree. For $i=1,\ldots,b-1$ let $v_i$ denote the left son of $v_{i-1}$. For every $i\in\{0,1,\ldots,b-1\}$ and $x\in\{0,1\}^{\ell-1}$ we define the random variables $a_i(x,k)=a_{v_i}(0^\ell,1\circ x,k)$, $h_i(x)=h_{v_i}(0^{\ell-1},x)$.
Define $\delta=c_1\epsilon\log^{-1}(1/\epsilon)$, $\ell_0=12\lceil\epsilon^{-1}\log(1/\epsilon)\rceil$, for some constant $c_1\in[0,1]$ to be set later on. Observe that for every fixing of the sequence $\{R_i\}$, $T$ is a palette-alternating tree code with distance $\delta$ if and only if for every $x\in\{0,1\}^{\ell-1}$ and $i\in\{0,1,\ldots,b-1\}$ it holds that $h_i(x)\ge\delta\ell$. Indeed, by definition, $T$ is a palette-alternating tree code with distance $\delta$ if and only if for every vertex $v$, $\ell\ge1$, and every distinct $x,y\in\{0,1\}^{\ell-1}$ it holds that $h_v(x,y)\ge\delta\ell$. However, by Claim18, the random variable $h_v(x,y)$ is the same as the random variable $h_i(y-x)$ for $i=\operatorname{depth}(v)\bmod b$.
For $x\in\{0,1\}^{\ell-1}$ and $i\in\{0,1,\ldots,b-1\}$ denote by $E_i(x)$ the event $h_i(x)<\delta\ell$. Note that as $R_1=1$ and since $\operatorname{Tr}(1)=1$ we have that $h_i(x)\ge1$ for every $x$. Thus, for $|x|<\ell_0$ we have that $h(x)/(|x|+1)\ge1/\ell_0$. Therefore, in order to prove Theorem15 it suffices to prove that $\Pr[\bigcup_{|x|\ge\ell_0}\bigcup_{i=0}^{b-1}E_i(x)]<1$.
Indeed, this will give a bound of $\min(1/\ell_0,\delta)=\Omega(\epsilon\log^{-1}(1/\epsilon))$ on the distance.
Fix $x\in\{0,1\}^{\ell-1}$ and $i\in\{0,1,\ldots,b-1\}$. Observe that $E_i(x)$ holds if and only if there exists a set $\mathcal T\subseteq[\ell]$ of size $\lceil(1-\delta)\ell\rceil$ such that for every $k\in\mathcal T$, $a_i(x,k)=0$. For ease of readability we ignore the ceiling in the calculations below. Recall that $a_i(x,1),\ldots,a_i(x,\ell)$ are independent. Further, $1-1/b$ fraction of them are uniform over $\mathbb F_2$ whereas the remaining $1/b$ fraction are uniform over $\mathbb F_4$. Note that by our choice of parameters, $\delta<1/b$. Thus, for any $\gamma\ge0$ and a fixed $\mathcal T$, we have that
\begin{align*}
\Pr[\forall k\in\mathcal T\ a_i(x,k)=0]
&\le2^{-(1-1/b-\gamma)\ell}4^{-(1/b-\delta+\gamma)\ell}\\
&\le2^{-(1-1/b)\ell}4^{-(1/b-\delta)\ell}
=2^{-(1+1/b-2\delta)\ell}.
\end{align*}
By taking the union bound over the choice of $\mathcal T$, and using Lemma7, we get that
\[
\Pr[E_i(x)]\le\binom\ell{\lceil(1-\delta)\ell\rceil}2^{-(1+1/b-2\delta)\ell}
\le2^{-(1+1/b-2\delta-H(\delta))\ell}.
\]
By the union bound,
\begin{align}
\Pr[\bigcup_{|x|\ge\ell_0}\bigcup_{i=0}^{b-1}E_i(x)]
&\le\sum_{|x|\ge\ell_0}\sum_{i=0}^{b-1}\Pr[E_i(x)]\tag{5}\\
&\le b\cdot\sum_{\ell=\ell_0}^\infty2^{\ell-1}\cdot2^{-(1+1/b-2\delta-H(\delta))\ell}\notag\\
&=\frac b2\sum_{\ell=\ell_0}^\infty2^{(H(\delta)+2\delta-1/b)\ell}.\notag
\end{align}
By taking $c_1$ sufficiently small and using Lemma8, we get that $H(\delta)+2\delta-1/b\le-\epsilon/3$. Therefore, Equation(5) is bounded above by
\[
b\sum_{\ell=\ell_0}^\infty2^{-\epsilon\ell/3}
=b\cdot\frac{2^{-\epsilon\ell_0/3}}{1-2^{-\epsilon/3}}
\le\frac{b\epsilon^4}{1-2^{-\epsilon/3}}
\le\frac{2\epsilon^3}{1-2^{-\epsilon/3}},
\]
where the penultimate inequality follows by our choice of $\ell_0$ and the last inequality follows since $b=\lceil1/\epsilon\rceil$. One can verify that the above is strictly bounded by 1 for any $\epsilon<1/3$.
\end{proof}

\paragraph{Current formal application context (Section1.4)}
Based on palette-alternating tree codes, we devise a deterministic interactive coding scheme against adversarial errors that approaches capacity and thus matches the rate obtained by [13]. The advantage of our coding scheme is that given an explicit construction of palette-alternating tree codes, our scheme is efficient. We believe that the recent progress on tree code constructions [11,24] may eventually lead to constructions of palette-alternating tree codes. The coding scheme suggested in [13], on the other hand, relies on a certain counting argument, and it is not clear to us how to obtain an efficient scheme based on these ideas.
\begin{authorclaim}{Theorem6 (application context only)}
Let $\epsilon>0$. Assume there exists an explicit palette-alternating tree code with rate $1-\epsilon$ and distance $\delta=\Omega(\epsilon\cdot\log^{-1}(1/\epsilon))$ (which, computational aspects aside, we know exists by Theorem4). Then, there exists an efficient deterministic coding scheme against $\epsilon$-fraction of adversarial errors with rate $1-O(\sqrt{\epsilon\log(1/\epsilon)})$.
\end{authorclaim}

\begin{thebibliography}{32}
\bibitem{r1} S. Agrawal, R. Gelles, and A. Sahai. Adaptive protocols for interactive communication. In Proc. IEEE International Symposium on Information Theory (ISIT), pages 595–599, 2016.
\bibitem{r2} N. Alon, M. Braverman, K. Efremenko, R. Gelles, and B. Haeupler. Reliable communication over highly connected noisy networks. In Proc. ACM SIGACT-SIGOPS Symposium on Principles of Distributed Computing (PODC), pages 165–173, 2016.
\bibitem{r3} Z. Brakerski and Y. T. Kalai. Efficient interactive coding against adversarial noise. In Proc. IEEE Symposium on Foundations of Computer Science (FOCS), pages 160–166, 2012.
\bibitem{r4} Z. Brakerski, Y. T. Kalai, and M. Naor. Fast interactive coding against adversarial noise. Journal of the ACM (JACM), 61(6):35:1–30, 2014.
\bibitem{r5} Z. Brakerski and M. Naor. Fast algorithms for interactive coding. In Proc. ACM-SIAM Symposium on Discrete Algorithms (SODA), pages 443–456, 2013.
\bibitem{r6} M. Braverman. Towards deterministic tree code constructions. In Proc. ACM-SIGACT Innovations in Theoretical Computer Science Conference (ITCS), pages 161–167, 2012.
\bibitem{r7} M. Braverman and K. Efremenko. List and unique coding for interactive communication in the presence of adversarial noise. In Proc. IEEE Symposium on Foundations of Computer Science (FOCS), pages 236–245, 2014.
\bibitem{r8} M. Braverman, R. Gelles, J. Mao, and R. Ostrovsky. Coding for interactive communication correcting insertions and deletions. IEEE Transactions on Information Theory, 63(10):6256– 6270, 2017.
\bibitem{r9} M. Braverman and A. Rao. Towards coding for maximum errors in interactive communication. In Proc. ACM Symposium on Theory of Computing (STOC), pages 159–166, 2011.
\bibitem{r10} M. Braverman and A. Rao. Toward coding for maximum errors in interactive communication. IEEE Transactions on Information Theory, 60(11):7248–7255, 2014.
\bibitem{r11} G. Cohen, B. Haeupler, and L. J. Schulman. Explicit binary tree codes with polylogarithmic size alphabet. In Proceedings of the 50th Annual ACM SIGACT Symposium on Theory of Computing, pages 535–544. ACM, 2018.
\bibitem{r12} R. Gelles. Coding for interactive communication: A survey. Foundations and Trends in Theoretical Computer Science, 13(1-2):1–157, 2017. doi:10.1561/0400000079.
\bibitem{r13} R. Gelles, B. Haeupler, G. Kol, N. Ron-Zewi, and A. Wigderson. Towards optimal deterministic coding for interactive communication, 2016.
\bibitem{r14} R. Gelles, A. Moitra, and A. Sahai. Efficient coding for interactive communication. IEEE Transactions on Information Theory, 60(3):1899–1913, 2014.
\bibitem{r15} Ran Gelles, Ankur Moitra, and Amit Sahai. Efficient and explicit coding for interactive communication. In 2011 IEEE 52nd Annual Symposium on Foundations of Computer Science, pages 768–777. IEEE, 2011.
\bibitem{r16} M. Ghaffari and B. Haeupler. Optimal Error Rates for Interactive Coding II: Efficiency and List Decoding. In Proc. IEEE Symposium on Foundations of Computer Science (FOCS), pages 394–403, 2014.
\bibitem{r17} M. Ghaffari, B. Haeupler, and M. Sudan. Optimal error rates for interactive coding I: Adaptivity and other settings. In Proc. ACM Symposium on Theory of Computing (STOC), pages 794–803, 2014.
\bibitem{r18} B. Haeupler. Interactive Channel Capacity Revisited. In Proc. IEEE Symposium on Foundations of Computer Science (FOCS), pages 226–235, 2014.
\bibitem{r19} A. Jain, Y. T. Kalai, and A. B. Lewko. Interactive coding for multiparty protocols. In Proc. ACM Innovations in Theoretical Computer Science Conference (ITCS), pages 1–10, 2015.
\bibitem{r20} G. Kol and R. Raz. Interactive channel capacity. In STOC, volume 13, pages 715–724, 2013.
\bibitem{r21} E. Kushilevitz and N. Nisan. Communication complexity. In Advances in Computers, volume 44, pages 331–360. Elsevier, 1997.
\bibitem{r22} A. Lewko and E. Vitercik. Balancing communication for multi-party interactive coding. CoRR, abs/1503.06381, 2015. arXiv:1503.06381.
\bibitem{r23} C. Moore and L. J. Schulman. Tree codes and a conjecture on exponential sums. In Proc. ACMSIGACT Innovations in Theoretical Computer Science Conference (ITCS), pages 145–154, 2014.
\bibitem{r24} Anand Kumar Narayanan and Matthew Weidner. On decoding Cohen-Haeupler-Schulman tree codes. In Proceedings of the Fourteenth Annual ACM-SIAM Symposium on Discrete Algorithms, pages 1337–1356. SIAM, 2020.
\bibitem{r25} P. Pudlák. Linear tree codes and the problem of explicit constructions. Linear Algebra and its Applications, 490:124–144, 2016.
\bibitem{r26} S. Rajagopalan and L. J. Schulman. A coding theorem for distributed computation. In Proc. ACM Symposium on Theory of Computing (STOC), pages 790–799, 1994.
\bibitem{r27} A. Rao and A. Yehudayoff. Communication complexity (early draft), 2018.
\bibitem{r28} L. J. Schulman. Communication on noisy channels: a coding theorem for computation. Proc. IEEE Symposium on Foundations of Computer Science (FOCS), pages 724–733, 1992.
\bibitem{r29} L. J. Schulman. Deterministic coding for interactive communication. In Proc. ACM Symposium on Theory of Computing (STOC), pages 747–756, 1993.
\bibitem{r30} L. J. Schulman. Postscript of 21 September 2003 to Coding for Interactive Communication. http://users.cms.caltech.edu/\textasciitilde{}schulman/Papers/intercodingpostscript.txt, 1994.
\bibitem{r31} L. J. Schulman. Coding for interactive communication. IEEE Transactions on Information Theory, 42(6):1745–1756, 1996.
\bibitem{r32} A. A. Sherstov and P. Wu. Optimal interactive coding for insertions, deletions, and substitutions. In Proc. IEEE Symposium on Foundations of Computer Science (FOCS), 2017.
\end{thebibliography}
\end{document}
